\section{Methods}
\label{sec:methods}

\subsection{Data and constraints}

Both parts come from the HINT benchmark~\cite{fu2022hint} and ClinicalTrials.gov~\cite{ctgov}.
In Part 1, the 312 latest trials (2010--2016) form the test set and the remaining 1,249 are split at random into train (1,171) and validation (78).
Only 157 train trials are labeled: the earliest 117 plus 40 chosen for rare reasons.
Recruitment covers 64.3\% of them and Funding 5.1\% (8 positives); 147 have exactly one reason.
The input is a 150-dimensional reason embedding (a PCA-reduced, randomly rotated sentence embedding).
In Part 2 each row is one trial in one phase.
Train (8,037 rows) and validation (1,013) hold trials that started in 2013 or earlier, split by trial; the 1,500 test rows come from 2013--2022.
Train success rates are 0.566, 0.477 and 0.653 for phases I, II and III.

Reason embeddings and labels are used in Part 2 only for the error analysis, never as classifier inputs.
Trial IDs are used only to match the two parts, and no deep learning is used.
The Part 1 semi-supervised methods include the unlabeled test embeddings in the unlabeled pool, which makes the setting transductive; no test labels are used.
Validation trials and CV hold-out folds are never in the pool.

\subsection{Part 1: recovering reason labels}

\paragraph{Shared components.}
Embeddings are standardized with train-file statistics.
The unlabeled pool $X_U$ holds the 1,014 masked train trials and the 312 test trials.
Each reason gets an L2-regularized LR with \texttt{class\_weight="balanced"}.
Pseudo-labels are down-weighted by \texttt{sample\_weight}; scikit-learn multiplies the two weights, so a rare-class pseudo-label ends up with more weight than $w_u$.
Every trial has at least one reason and an empty prediction scores 0, so all methods keep labels with probability $\ge 0.5$ and fall back to the most probable label when none passes.

\paragraph{Evaluation.}
Changing the prediction for one of the 78 validation trials moves sample-F1 by about 1.3 points, so we also report 2$\times$5-fold cross-validation (CV) on the 235 labeled trials (157 train plus 78 validation), with hold-out labels hidden.

\paragraph{Cluster diagnostics.}
All 1,561 trials are L2-normalized and reduced to 50 principal components.
K-means is run for $K$ from 2 to 80, recording inertia, silhouette, the label-set purity of labeled trials, and the adjusted Rand index (ARI)~\cite{hubert1985comparing} between clusters and label-sets.

\paragraph{Labeling methods.}
Nine families, tuned on validation and CV:

\begin{itemize}[leftmargin=*,itemsep=1pt,topsep=2pt]
\item \textbf{A, supervised LR} on the 157 labeled trials, $C\in\{0.01,0.1,1,10\}$.
\item \textbf{B0 / B, K-means cluster-then-label.}
K-means on L2-normalized PCA-50 features, centroids initialized with k-means++~\cite{arthur2007kmeans}, 10 restarts keeping the lowest inertia.
Cluster label vectors are smoothed toward the labeled prior $\bar y$: $\hat y_k=(\sum_{i\in L,c_i=k}y_i+\bar y)/(n_k+1)$.
B0 predicts $\hat y_k$; B uses it as a pseudo-label with weight $w_u\in\{0.3,1\}$ and retrains LR. $K\in\{10,20,40,80\}$.
\item \textbf{C, K-means $\cap$ LR agreement.}
Unlabeled trials where cluster voting and the supervised LR agree on all five labels get the LR labels as pseudo-labels ($w_u=1$).
As in co-training, the two models make different kinds of errors. $K\in\{10,20,40\}$.
\item \textbf{D, GMM soft labels.}
Diagonal-covariance GMM on PCA-20, \texttt{reg\_covar=1e-3}, 3 restarts, components initialized from K-means (scikit-learn default).
Each unlabeled trial gets a posterior-weighted soft label $p$ and enters LR as a $y=1$ row with weight $w_u p$ and a $y=0$ row with weight $w_u(1-p)$, $w_u=0.5$.
\item \textbf{E0 / E, semi-supervised GMM} following Nigam~\etal~\cite{nigam2000text}.
One Gaussian component per label-set observed in the labeled data: 7 label-sets, so $K=7$ (Funding's 8 positives fall into three of them).
Components share one covariance with a ridge of $10^{-2}I$, since 157 trials cannot support per-component covariances.
EM runs in a $d$-dimensional PCA space.
\emph{Initialization:} the first M-step uses the labeled one-hot responsibilities only, with no random or K-means start; labeled responsibilities stay fixed.
\emph{E-step:} $r_{ik}\propto\pi_k\mathcal N(z_i\mid\mu_k,\Sigma)$ for unlabeled trials.
\emph{M-step:} unlabeled responsibilities are weighted by $\lambda$, so $\lambda=0$ is a purely supervised generative classifier.
\emph{Convergence:} relative change in unlabeled log-likelihood below $10^{-5}$, at most 100 iterations.
Predictions are $P(y_j=1\mid z)=\sum_k p(k\mid z)\,c_{kj}$, with $c_k$ the label-set of component $k$.
E0 uses this posterior directly ($d\in\{10,20,40\}$, $\lambda\in\{0,0.3,1\}$).
E keeps unlabeled trials with maximum posterior above 0.8 or 0.95 and uses their MAP label-set as a pseudo-label ($w_u=0.3$).
\item \textbf{F, self-training LR.} Up to 5 rounds; a trial joins when all five probabilities satisfy $|2p_j-1|\ge 0.8$ or $0.95$.
\item \textbf{G, label spreading}~\cite{zhou2004learning} on a kNN graph ($k\in\{10,20\}$, $\alpha=0.5$), per reason.
\item \textbf{H, ensemble:} the mean probability of A ($C=0.1$), C ($K=20$) and E ($d=20$, threshold 0.95), one discriminative model and two pseudo-label models of different kinds.
\end{itemize}

\paragraph{Features.}
On the supervised LR we added, one at a time, \texttt{start\_year}; TF-IDF of \texttt{criteria}, \texttt{diseases\_cleaned} and \texttt{drugs} reduced to 50 dimensions by SVD; a matched-drug flag with multi-hot drug actions; multi-hot level-0 Reactome pathways~\cite{gillespie2022reactome}; and all of these.
\texttt{pathway\_level} describes the pathway, not the drug, so it was used only to select top-level pathways.
The chosen model is refit on all 235 labeled trials.

\subsection{Part 2: predicting trial outcomes}

\paragraph{Baseline.}
One LR per phase on a binary bag of words over the four text fields ($C=1$, up to 20,000 features).

\paragraph{Representations and classifiers.}
Bag of words, TF-IDF on unigrams, and TF-IDF on 1--2-grams, each with LR ($C$ from 0.01 to 10), Multinomial NB on counts or TF-IDF ($\alpha$ from 0.1 to 10), or distance-weighted cosine kNN ($k\in\{25,75,150\}$).

\paragraph{Field-wise TF-IDF.}
In one concatenated string, the long \texttt{criteria} field outweighs \texttt{drugs} and \texttt{icdcodes}.
We fit one TF-IDF per field and concatenate them with fixed weights (1.0 for \texttt{criteria}, 0.5 for the others).
Each ICD code is preceded by its three-character category (\texttt{E11.65} becomes \texttt{E11 E11.65})~\cite{who_icd10}, so related codes share a feature.
Optional blocks add drug-knowledge tokens (targets and actions prefixed with \texttt{TGT\_} and \texttt{ACT\_}) and RDKit~\cite{rdkit} Morgan fingerprints~\cite{rogers2010ecfp} (radius 2, 1,024 bits, bitwise maximum over a trial's drugs, rows scaled by the square root of the number of set bits).

\paragraph{Pooled LR.}
The project allows a model that uses the phase.
We use the feature augmentation of Daum\'e III~\cite{daume2007frustratingly}, with each phase as a domain:
\begin{equation}
\phi(x,p)=\big[\,x,\ \mathbb 1_{p=\mathrm{I}}\,x,\ \mathbb 1_{p=\mathrm{II}}\,x,\ \mathbb 1_{p=\mathrm{III}}\,x,\ e_p\,\big].
\end{equation}
Phase $p$ has effective weights $w_0+w_p$ and its own intercept from $e_p$.
L2 regularization shrinks $w_p$ toward 0, so phases share strength.
One LR is fit on all rows, $C\in\{0.3,1,3\}$.
We also averaged it with a per-phase TF-IDF LR, kNN or MNB.

\paragraph{Error analysis.}
Failed validation trials are matched by \texttt{trial\_id} to Part 1 IDs and take the true or predicted reasons; unmatched trials form a ``no recorded reason'' group, and a trial can count toward several reasons.
Each phase is analyzed separately (we assert unique \texttt{trial\_id} within a phase).
For each group we report its size, how many reasons are true or predicted labels, and the fraction predicted successful, which is the group's false positive rate (FPR).

\paragraph{Trade-offs and drift.}
We tested thresholds from 0.4 to 0.7, \texttt{class\_weight="balanced"}, shrinkage toward the phase success rate $\pi$, $p'=\pi+s(p-\pi)$ with $s\in\{1,0.9,0.8,0.7\}$, and clipping to $[0.02,0.98]$, which caps the cost of a confident mistake at $-\ln 0.02\approx 3.9$.
As a proxy for the later test years, we train on trials from 2008 or earlier (6,326 rows) and evaluate on 2010 or later (2,290 rows).
The final model is refit on train $\cup$ validation (9,050 rows).
All models use scikit-learn~\cite{pedregosa2011scikit}, and every figure is produced by code in the submitted notebook.
